\documentclass[11pt]{scrartcl}
\usepackage{geometry}
\usepackage{hyperref}
\usepackage{xcolor}
\usepackage{minted}
\usepackage[sfdefault]{carlito}
\renewcommand{\familydefault}{\sfdefault}
\geometry{margin=1in}

\definecolor{xhbg}{HTML}{FAFAFA}
\setminted{
  fontsize=\small,
  breaklines=true,
  autogobble=true,
  frame=single,
  framesep=2mm,
  linenos=true,
  tabsize=2,
  bgcolor=xhbg
}
\usemintedstyle{friendly}

\newminted[code]{haskell}{}

\title{Xeus-Haskell: Manual and Literate Sources}
\author{Masaya Taniguchi and xeus-haskell contributors}
\date{}

\begin{document}
\maketitle
\tableofcontents
\newpage

\section{Welcome to Xeus-Haskell}


Xeus-Haskell provides Jupyter kernels for Haskell based on the native implementation of the Jupyter protocol, \href{https://github.com/jupyter-xeus/xeus}{xeus}.
\subsection{Preliminaries}


To start using or developing Xeus-Haskell, please note the following:
\begin{itemize}
\item \textbf{MicroHs Backend:} The \texttt{xhaskell-mhs} kernel uses \href{https://github.com/augustss/MicroHs}{MicroHs}. It runs natively and in WebAssembly without GMP; supported libraries and extensions follow MicroHs.
\item \textbf{GHC Backend:} The \texttt{xhaskell-ghc} kernel uses the native-bignum flavour (without GMP) of a pinned GHC Wasm toolchain and runs in JupyterLite.
\item \textbf{Xeus Protocol:} Both kernels implement the Jupyter protocol through \href{https://github.com/jupyter-xeus/xeus}{xeus}. The native build currently installs \texttt{xhaskell-mhs}.
\item \textbf{WebAssembly Support:} JupyterLite registers both \texttt{xhaskell-mhs} and \texttt{xhaskell-ghc}.
\end{itemize}


\subsection{Quickstart}


Since this project is not yet available on Conda-forge or Emscripten-forge, you need to build it from source. We have organized the workflows using \href{https://pixi.sh}{pixi}.
\subsubsection{For Native JupyterLab}


\begin{minted}{text}
git clone https://github.com/jupyter-xeus/xeus-haskell
cd xeus-haskell
pixi run -e native native
pixi run -e native serve # JupyterLab is ready!
\end{minted}


\subsubsection{For WebAssembly JupyterLite}


\begin{minted}{text}
git clone https://github.com/jupyter-xeus/xeus-haskell
cd xeus-haskell
pixi run -e wasm-build wasm
pixi run -e wasm-build serve # JupyterLite is ready!
\end{minted}

Each public Pixi environment composes the feature with the same name:
\begin{itemize}
\item \texttt{native}: native build and tests
\item \texttt{docs}: documentation generation
\item \texttt{wasm-build}: Emscripten and JupyterLite build tools
\item \texttt{wasm-host}: Emscripten target libraries
\item \texttt{browser-test}: Chromium integration tests
\end{itemize}
The historically named \texttt{wasm-host} environment is an
\texttt{emscripten-wasm32} target prefix, not another host toolchain; this is
why its installation specifies \texttt{--platform emscripten-wasm32}.

\subsection{A Deep Dive into Xeus-Haskell}


\begin{itemize}
\item \hyperref[sec:contributing]{Contributing Guide}
\item \hyperref[sec:communication]{Communication Protocol}
\item \hyperref[sec:repl]{REPL Mechanism}
\item \href{https://github.com/jupyter-xeus/xeus-haskell/blob/main/AGENTS.md}{AI Agent Guide}
\end{itemize}

\section{Communication and System Architecture}
\label{sec:communication}


This document describes the common Jupyter and xeus layers, followed by the MicroHs and GHC execution paths.
\subsection{System Stack}


The communication stack is organized as follows:
Jupyter Frontend \(\leftrightarrow\) xeus/xeus-lite \(\leftrightarrow\) backend interpreter \(\leftrightarrow\) MicroHs runtime or persistent GHC Wasm session.


\subsection{Layer Details}


\subsubsection{Jupyter and Xeus Protocol Layer}


This layer handles the ZeroMQ-based Jupyter protocol. It manages the heartbeat, shell, and control sockets.
\subsubsection{Xeus to Kernel Interpreter Layer}


The \texttt{xinterpreter} implementation hooks into the Xeus kernel lifecycle. It receives execution requests and delegates them to the Haskell backend.
\subsubsection{MicroHs Bridge to Haskell Runtime}


This is the bridge layer between C++ and the Haskell-compiled Compiler API wrapper (\texttt{Repl.lhs}).
The \texttt{mhs\_repl.cpp} bridge performs three primary tasks:
\begin{enumerate}
\item \textbf{Initialize MicroHs Runtime}:
   The \texttt{MicroHsRepl} constructor locates runtime files (standard libraries) using \texttt{microhs\_runtime\_dir}. After initialization, it evaluates a trivial expression (\texttt{"0"}) to warm up the compiler and cache library compilation.
\item \textbf{Capture Standard Output}:
   To receive output from the Haskell world, we use POSIX pipes to redirect \texttt{stdout}.
   \begin{itemize}
   \item C++ sets up a redirection from \texttt{stdout} to a temporary buffer.
   \item After Haskell evaluation completes, the buffer is read.
   \item For implementation details, see \texttt{capture\_stdout} in the source code.
   \end{itemize}
\item \textbf{Parse Captured Content}:
   The output is parsed to distinguish plain text from rich content (HTML, LaTeX, etc.) using the internal protocol described below.
\end{enumerate}


\subsubsection{REPL Engine and MicroHs Compiler Layer}


This layer uses the MicroHs Compiler API to evaluate Haskell code strings in the runtime environment.

\subsection{GHC Browser Kernel}
The browser-only \texttt{xhaskell-ghc} kernel separates protocol handling from compiler execution:
\begin{enumerate}
\item The kernelspec publishes the compressed GHC root filesystem and dynamic-linker modules as shared JupyterLite assets.
\item \texttt{browser/ghc\_support.js} creates an in-memory WASI filesystem, extracts \texttt{rootfs.tar.zst}, loads \texttt{libxeus-haskell-ghc.so}, and initializes the exported GHC API entry point.
\item \texttt{haskell/Playground.hs} retains one interpreter-backed GHC session and dispatches imports, declarations, statements, expressions, and supported GHCi colon commands through the GHC API.
\item An Emscripten run dependency evaluates \texttt{1 + 1} during startup, then restores the clean interactive state so the warm-up does not leak \texttt{it}.
\item Cells follow GHCi semantics and keep declarations, imports, IO state, and \texttt{it} in the retained session. A declaration prefix may be followed by an expression or executable statements; the declarations remain available if the suffix fails. Completion, inspection, and completeness checks query that same state.
\item The C++ interpreter uses Emscripten JSPI to await serialized JavaScript requests and translates their JSON payloads into Jupyter replies.
\end{enumerate}

\subsection{Rich Content Protocol}


Both \texttt{xhaskell-mhs} and \texttt{xhaskell-ghc} support rich HTML, LaTeX, and Markdown by embedding control sequences in the standard output stream. The kernel decodes each valid rich frame into a Jupyter \texttt{display\_data} message; plain text remains an execution result or stream according to its source.
\subsubsection{Protocol Structure}


Rich content is strictly delimited by \textbf{ASCII control characters}:
\begin{minted}{text}
<STX><MIME_TYPE><US><CONTENT><ETX>
\end{minted}


\begin{center}
\begin{tabular}{|l|c|p{8cm}|}
\hline
\textbf{Component} & \textbf{ASCII (Hex)} & \textbf{Description} \\
\hline
\texttt{STX} & \texttt{02} & \textbf{S}tart of \textbf{T}e\textbf{x}t \\
\hline
\texttt{MIME\_TYPE} & - & The target MIME type (e.g., \texttt{text/html}) \\
\hline
\texttt{US} & \texttt{1F} & \textbf{U}nit \textbf{S}eparator \\
\hline
\texttt{CONTENT} & - & The raw data/string to be rendered \\
\hline
\texttt{ETX} & \texttt{03} & \textbf{E}nd of \textbf{T}e\textbf{x}t \\
\hline
\end{tabular}
\end{center}
\subsubsection{Supported Content Types}


\begin{itemize}
\item \texttt{text/plain}: Standard console output.
\item \texttt{text/html}: Rendered as HTML in Jupyter.
\item \texttt{text/latex}: Mathematical equations rendered via MathJax.
\item \texttt{text/markdown}: Formatted text and tables.
\end{itemize}


All content is expected to be \textbf{UTF-8} encoded.
\subsection{Haskell Display API}


Both kernels expose the protocol through the compatible \texttt{XHaskell.Display} module and its \texttt{Display} typeclass.
\subsubsection{Display Typeclass Interface}


To enable rich output for a custom type, implement the \texttt{Display} typeclass:
\begin{minted}{haskell}
import XHaskell.Display

-- | Data structure for the protocol
data DisplayData = DisplayData {
    mimeType :: String,
    content  :: String
}
\end{minted}


\subsubsection{Implementation Examples}


\begin{minted}{haskell}
newtype HTMLString = HTMLString String
instance Display HTMLString where
    display (HTMLString s) = DisplayData "text/html" s

newtype LaTeXString = LaTeXString String
instance Display LaTeXString where
    display (LaTeXString s) = DisplayData "text/latex" s

-- Default instance for plain strings
instance Display String where
    display s = DisplayData "text/plain" s
\end{minted}


\subsubsection{Protocol Generation}


The \texttt{DisplayData} record implements a \texttt{Show} instance that automatically formats the string into the \texttt{<STX>...<ETX>} protocol.
\begin{minted}{haskell}
putStr $ show (display (HTMLString "<p style='color: red'>Hello Haskell!</p>"))
\end{minted}


\texttt{02} \texttt{text/html} \texttt{1F} \texttt{<p style='color: red'>Hello Haskell!</p>} \texttt{03}

\section{Contributing to xeus-haskell}
\label{sec:contributing}


Thank you for your interest in contributing to \texttt{xeus-haskell}!
Xeus and xeus-haskell are subprojects of Project Jupyter and are subject to the \href{https://github.com/jupyter/governance}{Jupyter governance} and \href{https://github.com/jupyter/governance/blob/master/conduct/code\_of\_conduct.md}{Code of conduct}.
\subsection{General Guidelines}


For general documentation about contributing to Jupyter projects, see the \href{https://jupyter.readthedocs.io/en/latest/contributor/content-contributor.html}{Project Jupyter Contributor Documentation}.
\subsection{Community}


The Xeus team organizes public video meetings. The schedule for future meetings and minutes of past meetings can be found on our \href{https://jupyter-xeus.github.io/}{Team Compass}.
\subsection{Getting Started}


\subsubsection{Prerequisites}


The project uses \href{https://pixi.sh/}{Pixi} for dependency management and workflow automation. If you haven't installed it yet:
\begin{minted}{text}
curl -fsSL https://pixi.sh/install.sh | sh
\end{minted}


\subsubsection{Setting Up Your Environment}


\begin{enumerate}
\item Fork the repository on GitHub.
\item Clone your fork locally:
\end{enumerate}


\begin{minted}{text}
git clone https://github.com/<your-username>/xeus-haskell.git
cd xeus-haskell
\end{minted}


\begin{enumerate}
\item Initialize the development environment:
\end{enumerate}


\begin{minted}{text}
pixi run -e native native
\end{minted}


The \texttt{native} environment includes all necessary tools (CMake, C++ compilers, Python for testing, etc.).
\subsection{Development Workflow}


\subsubsection{Building the Kernel}


To build and install the kernel in your local environment:
\begin{minted}{text}
pixi run -e native native
\end{minted}


\subsubsection{Running Tests}


Tests are organized into three deterministic tiers with a shared quality
contract. Build and install the WebAssembly kernels before running Tier 3:
\begin{minted}{text}
# Native C++ and Python contracts
pixi run -e native test

# Production Wasm kernels in Chromium
pixi run -e wasm-build test
\end{minted}

The three tiers share execution, mixed-cell, persistence, error, completion,
inspection, completeness, silent-output, and rich-display quality concepts.
The Playwright tier applies them to the production MicroHs and GHC kernels.
MicroHs receives the same library mounts as the JupyterLite tasks. GHC uses its
production root filesystem, pinned WASI shim, and generated dynamic linker.
HTTP caching is disabled for the static test server, so each GHC run fetches
and extracts \texttt{rootfs.tar.zst}. Chromium executes the actual
Web Worker and \texttt{importScripts} boundary used by the browser kernels.


\subsubsection{Working with WebAssembly (JupyterLite)}


If you are contributing to the WebAssembly build, you need the \texttt{wasm-build} and \texttt{wasm-host} environments:
\begin{minted}{text}
# Download inputs and build both Wasm kernels
pixi run -e wasm-build wasm

# Build and audit the distributable JupyterLite site
pixi run -e wasm-build pages

# Run the JupyterLite server locally
pixi run -e wasm-build serve
\end{minted}

Runtime notices are installed below \texttt{share/licenses/xeus-haskell}.
The JupyterLite build publishes that tree through an Xeus mount at
\texttt{/usr/share/licenses/xeus-haskell}. It includes generated notices, \texttt{licenses.toml},
and a \texttt{SOURCE\_OFFER.md} file identifying the verified corresponding
source bundle for that deployment. The bundle is published beside the static
JupyterLite site and as an asset of tagged GitHub releases.


\subsection{Project Structure}


\begin{itemize}
\item \texttt{cmake/}: Reusable CMake modules and package configuration templates.
\item \texttt{docs/}: The LaTeX manual, documentation assets, and generated PDF directory.
\item \texttt{notebooks/}: Example notebooks for the assembled kernels.
\item \texttt{test/cpp/}: Tier 1 native runtime and portable protocol adapter tests.
\item \texttt{test/python/}: Tier 2 Jupyter protocol and canonical GHC reply tests.
\item \texttt{test/browser/}: Tier 3 Chromium tests for both production Wasm kernels.
\item \texttt{xhaskell/common/}: Headers and browser glue shared by both kernels.
\item \texttt{xhaskell/microhs/}: The native and WebAssembly MicroHs kernel.
\item \texttt{xhaskell/ghc/}: The browser-only GHC kernel and resource toolchain.
\item Backend directories use \texttt{browser/}, \texttt{haskell/},
  \texttt{include/}, \texttt{src/}, and \texttt{share/} by responsibility.
\item Generated state belongs in the top-level build directories or the
  ignored local GHC toolchain and resource directories.
\end{itemize}


\subsection{Style Guidelines}


\begin{itemize}
\item \textbf{C++}: Follow modern C++ practices. We use Xeus as the base library.
\item \textbf{Haskell}: Keep MicroHs compatibility constraints in the MicroHs backend; GHC-specific APIs and extensions belong in the GHC backend.
\item \textbf{Documentation}: Write documentation in LaTeX (\texttt{.tex}) and keep examples runnable with the current Pixi workflows.
\end{itemize}

\subsubsection{Building Documentation}

Generate the project PDF documentation with:
\begin{minted}{text}
pixi run -e docs docs
\end{minted}

The generated manual is written to \texttt{docs/\_build/xhaskell.pdf}.



\subsection{Submitting Changes}


\begin{enumerate}
\item Create a new branch for your feature or bugfix.
\item Ensure all tests pass.
\item Submit a Pull Request with a clear description of the changes.
\end{enumerate}


\clearpage
\section{MicroHs Literate Haskell Source}
\input{../xhaskell/microhs/haskell/Repl.lhs}
\subsection{Snippet Analysis and Parsing Module}
\input{../xhaskell/microhs/haskell/Repl/Analysis.lhs}
\subsection{Compilation and Execution Translation Module}
\input{../xhaskell/microhs/haskell/Repl/Compiler.lhs}
\subsection{Persistent Context and Definition Storage Module}
\input{../xhaskell/microhs/haskell/Repl/Context.lhs}
\subsection{Error Classification and Rendering Module}
\input{../xhaskell/microhs/haskell/Repl/Error.lhs}
\subsection{Runtime Command Execution Module}
\input{../xhaskell/microhs/haskell/Repl/Executor.lhs}
\subsection{Shared Utility Functions Module}
\input{../xhaskell/microhs/haskell/Repl/Utils.lhs}

\end{document}
